\documentclass[11pt]{article}

\usepackage[margin=1in]{geometry}
\usepackage{amsmath}
\usepackage{booktabs}
\usepackage{enumitem}
\usepackage{hyperref}

\title{CTMC Submission Pipeline}
\author{M148 Capstone Project}
\date{\today}

\begin{document}
\maketitle

\section{Goal}

The Kaggle task requires one probability per open journey:
\[
\Pr(\texttt{order\_shipped}=1 \mid \text{observed journey so far}).
\]
The required submission format is
\begin{verbatim}
id,order_shipped
-1000001271 551641434,0.1
-100001164 -1710062169,0.95
\end{verbatim}

The CTMC pipeline writes this exact two-column format.  The standalone script
writes under \texttt{results/submissions/} by default, while the exploration
notebook writes run outputs under \texttt{results/}.  The script is
\texttt{src/models/ctmc\_submission.py}.

\section{Training Data}

The CTMC models use the checkpointed event data produced by
\texttt{src/data\_engineering/}.  In particular, the main input is
\texttt{data/journeys\_flattened.parquet}, where each journey has an ordered
list of pairs
\[
(t_m, s_m),
\]
with timestamp \(t_m\) and action state \(s_m=\texttt{ed\_id}\).

From each journey, the CTMC pipeline constructs observed transitions
\[
(i,j,\Delta t)=(s_m,s_{m+1},t_{m+1}-t_m).
\]
The success state is
\[
s^\star = 28,
\]
which corresponds to \texttt{order\_shipped}.

\section{Method 1: Global CTMC}

The global model estimates a single generator matrix \(Q\) from all training
journeys.  For each transition \(i\to j\), define
\[
N_{ij}=\#\{m:s_m=i, s_{m+1}=j\}
\]
and
\[
T_i=\sum_{m:s_m=i}(t_{m+1}-t_m).
\]
The off-diagonal transition rates are estimated as
\[
q_{ij}=\frac{N_{ij}}{T_i}, \qquad i\neq j.
\]
The diagonal is set so each row sums to zero:
\[
q_{ii}=-\sum_{j\neq i}q_{ij}.
\]

For prediction, \texttt{order\_shipped} is treated as an absorbing success
state by setting its row in \(Q\) to zero.  Given current state \(s\), the
submission probability is
\[
\Pr(\text{success by }60\text{ days}\mid s)
= \left[\exp(Q_{\mathrm{abs}}\,h)\right]_{s,s^\star},
\]
where \(h=60\cdot24\cdot60\cdot60\) seconds.

\section{Method 2: Clustered CTMC}

The segmented model first clusters journeys without using the purchase label.
Features include:

\begin{itemize}[leftmargin=*]
\item action counts and proportions,
\item total observed time and average inter-action gap,
\item time spent in common states,
\item first and current action,
\item time to key actions when observed.
\end{itemize}

After clustering, the pipeline fits one generator matrix per segment:
\[
Q^{(k)}, \qquad k=1,\ldots,K.
\]
For an open test journey, the same leakage-safe features are computed, the
cluster is predicted, and the probability is
\[
\left[\exp(Q^{(k)}_{\mathrm{abs}}h)\right]_{s,s^\star}.
\]

\section{Method 3: Neural Exponential-Rate CTMC}

The neural model estimates a customer-dependent success intensity directly.  It
does not predict the next state.  Given journey features \(x\), it learns
\(\lambda(x)\) and uses the exponential waiting-time model
\[
\Pr(T_{\mathrm{success}}\le t\mid x)=1-\exp(-\lambda(x)t).
\]

Successful training snapshots use the observed remaining time to
\texttt{order\_shipped} when available.  Incomplete snapshots are treated as
censored at the 60-day horizon with a small baseline hazard.  The resulting
probability is written directly in the Kaggle format.

\section{Outputs}

When \texttt{data/open\_journeys1.csv} is present, run:
\begin{verbatim}
python src/models/ctmc_submission.py
\end{verbatim}

The standalone script writes:
\begin{itemize}[leftmargin=*]
\item \texttt{results/submissions/ctmc\_global\_submission.csv}
\item \texttt{results/submissions/ctmc\_clustered\_submission.csv}
\item \texttt{results/submissions/ctmc\_neural\_rate\_submission.csv}
\item \texttt{results/submissions/ctmc\_blend\_submission.csv}
\end{itemize}

The notebook writes the same CTMC files to \texttt{results/}, along with
tabular Random Forest and XGBoost submissions.

If \texttt{data/open\_journeys1\_flattened\_all0.csv} is available, it is used
only to match the official Kaggle row order.  If it is missing but
\texttt{data/open\_journeys1.csv} is present, the script creates an all-zero
template with the required \texttt{id,order\_shipped} columns from the unique
open journey IDs.

\section{Comparison to Existing Random Forest}

The existing Random Forest submission path trains on truncated journey
features and writes the same columns:
\[
\texttt{id},\quad \texttt{order\_shipped}.
\]
The CTMC notebook now includes a comparison cell that plots CTMC variants
against logistic regression, gradient boosting, and random forest baselines.
The official Kaggle comparison should still be made with held-out or submitted
test scores, because notebook scores are internal validation metrics.

\section{Important Caveat}

The flattened all-zero file is not a feature table; it is an output template.
The CTMC features come from the open journey event history.  Therefore the CTMC
submission script needs \texttt{open\_journeys1.csv}, or an equivalent event file
with at least:
\[
\texttt{id},\quad \texttt{event\_timestamp},\quad \texttt{ed\_id}.
\]

\end{document}
